% Source: https://artofproblemsolving.com/community/c3h2184819_triangle_with_integer_length_sides
% Stable identity: aops-topic-2184819-q01
Let the two equal sides be $a$ and the base $b$, and put $x=\dfrac {b}{2}$ and $y=a-x$. The incircle touches the base at its midpoint, so the tangent lengths are $x$ from each end of the base and $y$ from the apex. The semiperimeter is $s=a+x=2x+y$, with $s-a=x$ and $s-b=y$.

\begin{asy}
usepackage("amsmath");
size(5.5cm);
defaultpen(fontsize(10pt));
pair B=(-24,0), C=(24,0), A=(0,32), I=(0,12);
draw(A--B--C--cycle,linewidth(1));
draw(circle(I,12),blue);
pair T1=(0,0), T2=B+24*unit(A-B), T3=C+24*unit(A-C);
draw(I--T1,red+dashed); draw(I--T2,red+dashed);
dot(I); dot(T1); dot(T2); dot(T3);
label("$A$",A,N); label("$B$",B,SW); label("$C$",C,SE); label("$I$",I,E);
label("$r=12$",(I+T1)/2,E,red+fontsize(8pt));
label("$x$",(B+T1)/2,S,fontsize(9pt)); label("$x$",(T1+C)/2,S,fontsize(9pt));
label("$x$",(B+T2)/2,NW,fontsize(9pt)); label("$y$",(T2+A)/2,NW,fontsize(9pt));
label("$x$",(C+T3)/2,NE,fontsize(9pt)); label("$y$",(T3+A)/2,NE,fontsize(9pt));
label("\mbox{}",(min(currentpicture,user=true)+max(currentpicture,user=true))/2,black); // keeps the text after the figure black when LaTeX typesets the labels inline
\end{asy}

\textbf{The inradius is $12$.} The area is $K=rs$ and the perimeter is $2s$, so $K=6\cdot2s$ means $r=12$.

\textbf{An equation for $x$ and $y$.} Heron's formula gives $K^2=s(s-a)^2(s-b)=s\,x^2y$. With $K=12s$ this becomes $144s=x^2y$, that is,
\[144(2x+y)=x^2y .\]
As a quadratic in $x>0$, $yx^2-288x-144y=0$, so
\[x=\dfrac{144+12\sqrt{144+y^2}}{y}.\]

\textbf{Integrality.} Here $2x=b$ and $2y=2a-b$ are integers, so $\sqrt{144+y^2}=\dfrac{xy-144}{12}$ is rational. With $Y=2y$, the number $\sqrt{576+Y^2}=2\sqrt{144+y^2}$ is the rational square root of an integer, hence an integer $Z$, and
\[(Z-Y)(Z+Y)=576 .\]
The two factors have the same parity and an even product, so both are even: $Z-Y=2u$, $Z+Y=2v$ with $uv=144$ and $u<v$. Then $y=\dfrac{v-u}{2}$, $\sqrt{144+y^2}=\dfrac{u+v}{2}$ and
\[x=\dfrac{12(u+v+24)}{v-u},\qquad b=2x,\qquad a=x+y .\]
\begin{center}
\begin{tabular}{c|ccccccc}
$(u,v)$ & $(1,144)$ & $(2,72)$ & $(3,48)$ & $(4,36)$ & $(6,24)$ & $(8,18)$ & $(9,16)$\\\hline
$y$ & $71.5$ & $35$ & $22.5$ & $16$ & $9$ & $5$ & $3.5$\\
$x$ & $\dfrac{156}{11}$ & $16.8$ & $20$ & $24$ & $36$ & $60$ & $84$\\
$(a,a,b)$ & -- & -- & $(42.5,42.5,40)$ & $(40,40,48)$ & $(45,45,72)$ & $(65,65,120)$ & $(87.5,87.5,168)$
\end{tabular}
\end{center}
In the first two columns $b=2x$ is not an integer, and in the third and last columns $a$ is not. So there are exactly three triangles:
\[\boxed{(40,40,48),\quad(45,45,72),\quad(65,65,120)}.\]
Check, for $(40,40,48)$: the height is $\sqrt{40^2-24^2}=32$, the area is $\dfrac{1}{2}\cdot48\cdot32=768$ and the perimeter is $128$; indeed $768=6\cdot128$. Similarly $(45,45,72)$ has area $972=6\cdot162$ and $(65,65,120)$ has area $1500=6\cdot250$.
